\uuid{11dq}
\titre{Résolution d'équation différentielle du premier ordre}
\auteur{bodin ; croizat ; sacre}
\organisation{exo7}
\licence{CC BY-NC-SA 4.0}
\datecreate{2022-03-08}
\niveau{L1}
\module{Analyse}
\chapitre{Equation différentielle}
\sousChapitre{Résolution d'équation différentielle du premier ordre}
\difficulte{2}
\theme{}
\exo7qcmid{hilisit.equations-differentielles.047}

\contenu{
\begin{question}
\qclassification{225.01}
Soit l'équation différentielle $y' = y + x^2-1$.
Quelles sont les affirmations vraies ?
\begin{answers}
  \bad{Les solutions de l'équation homogène associée sont $y(x) = \frac13x^3-x+C$.}
  \bad{Les solutions de l'équation homogène associée sont $y(x) = Ce^{x^2-1}$.}
  \bad{Une solution particulière est $y(x) = e^{x}$.}
  \good{Une solution particulière est $y(x) = -x^2-2x-1$.}
\end{answers}
\begin{explanations}
L'équation homogène est $y' = y$, dont les solutions sont les $y_h(x) = Ce^{x}$.
Une solution particulière de l'équation $y' = y + x^2-1$ est $y_p(x) = -x^2-2x-1$.
Les solutions générales sont alors $y(x) = y_h(x) + y_p(x)$.
\end{explanations}
\end{question}
}
